\documentclass[UTF8]{ctexart}

\usepackage{amscd,amsthm,amsfonts,amssymb,amsmath}
\usepackage[colorlinks=true]{hyperref}
\usepackage[all]{xy}
\usepackage{mathrsfs}
\usepackage{tikz}
\usepackage{bm}
\usepackage{geometry}
\usepackage{xcolor}
\usepackage{eso-pic}
% \usepackage{CJK}
\usepackage{arydshln}
% \usepackage[11pt]{moresize}



\newtheorem{thm}{定理}
\newtheorem{cor}[thm]{推论}
\newtheorem{lem}[thm]{引理}
\newtheorem{prop}[thm]{命题}
\newtheorem{defn}[thm]{定义}
\newtheorem{ques}[thm]{问题}
\newtheorem{eg}[thm]{例题}
\newtheorem{ex}[thm]{习题}
\newtheorem{rmk}[thm]{注释}
\newtheorem{ntn}[thm]{记号}




\begin{document}


\title{习题课材料（五）}
\date{}

\maketitle

%\noindent{\Large\textbf{注: 带 $\heartsuit $号的习题有一定的难度、比较耗时, 请量力为之.}}

\begin{ex}
  给定 \(\boldsymbol{a} \in \mathbb{R}^3\).
  计算 \(\boldsymbol{a}\) 与坐标向量 \(\boldsymbol{e}_1,\boldsymbol{e}_2,\boldsymbol{e}_3\)
  的夹角的余弦, 并计算这三个余弦值的平方和.
\end{ex}



\begin{ex}
  证明, 分块上三角矩阵
  \(X=\begin{bmatrix} c & \boldsymbol{a}^{\mathrm{T}} \\ \boldsymbol{0} & Q\end{bmatrix}\)
  是正交矩阵时, 必有 \(c = \pm 1\), \(\boldsymbol{a} = \boldsymbol{0}\),
  \(Q\) 是正交矩阵.
\end{ex}



\begin{ex}[\(\heartsuit\)]
  设可逆矩阵 \(A\) 保持向量之间角度不变.
  证明 \(A\) 是正交矩阵的常数倍.
\end{ex}



\begin{ex}
 设 \(f\colon \mathbb{R}^n \to \mathbb{R}\)
  为线性映射. 证明存在唯一 \(\boldsymbol{a} \in \mathbb{R}^n\),
  使得 \(f(\boldsymbol{x}) = \boldsymbol{a}^{\mathrm{T}}\boldsymbol{x}\).
\end{ex}



\begin{ex}
  设 \(\boldsymbol{a}_1, \boldsymbol{a}_2, \boldsymbol{a}_3\)
  为 \(\mathbb{R}^3\) 的标准正交基.
  证明
  \begin{equation*}
    \boldsymbol{b}_1 = \frac{1}{3}(2\boldsymbol{a}_1 + 2\boldsymbol{a}_2 - \boldsymbol{a}_3), \quad
    \boldsymbol{b}_2 = \frac{1}{3}(2\boldsymbol{a}_1 - 1\boldsymbol{a}_2 + 2\boldsymbol{a}_3), \quad
    \boldsymbol{b}_1 = \frac{1}{3}(\boldsymbol{a}_1 - 2\boldsymbol{a}_2 - 2\boldsymbol{a}_3)
  \end{equation*}
  也是 \(\mathbb{R}^3\) 的标准正交基.
\end{ex}


\begin{ex}
  向量 \(\boldsymbol{a}, \boldsymbol{b} \in \mathbb{R}^n\)
  围成一个三角形. 证明它的面积的平方为
  \(\frac{1}{4}(\|\boldsymbol{a}\|^{2} \cdot \|\boldsymbol{b}\|^2 - (\boldsymbol{a}^{\mathrm{T}}\boldsymbol{b})^2)\).
\end{ex}

\begin{ex}
  对向量组
  \begin{equation*}
    \boldsymbol{v}_1 =
    \begin{bmatrix}
      1 \\ 2 \\ -2
    \end{bmatrix},
    \boldsymbol{v}_2 =
    \begin{bmatrix}
      4 \\ 4 \\-3
    \end{bmatrix},
    \boldsymbol{v}_3 =
    \begin{bmatrix}
      -1 \\ 7 \\ 2
    \end{bmatrix}
  \end{equation*}
  进行 Schmidt 正交化.
\end{ex}


\begin{ex}
  求矩阵 \(\begin{bmatrix} 2 & 1 & -1 \\ 1 & 0 & 2 \\ 2 & -1 & 3 \end{bmatrix}\) 的 QR 分解.
\end{ex}
\end{document}
